\documentclass[11pt,a4paper]{article}

\usepackage{amsmath,amssymb,amsthm}
\usepackage{hyperref}
\usepackage[margin=1in]{geometry}
\usepackage{xcolor}
\usepackage{booktabs}
\usepackage{array}
\usepackage{float}
\usepackage{mathtools}

\newtheorem{theorem}{Theorem}[section]
\newtheorem{conjecture}[theorem]{Conjecture}
\newtheorem{corollary}[theorem]{Corollary}
\newtheorem{lemma}[theorem]{Lemma}
\newtheorem{proposition}[theorem]{Proposition}
\newtheorem{definition}[theorem]{Definition}
\theoremstyle{remark}
\newtheorem{remark}[theorem]{Remark}

\newcommand{\PT}{\mathbb{PT}}
\newcommand{\CP}{\mathbb{CP}}
\newcommand{\C}{\mathbb{C}}
\newcommand{\Z}{\mathbb{Z}}
\newcommand{\R}{\mathbb{R}}
\newcommand{\F}{\mathbb{F}}
\newcommand{\cO}{\mathcal{O}}
\newcommand{\cC}{\mathcal{C}}
\newcommand{\cN}{\mathcal{N}}
\newcommand{\cU}{\mathcal{U}}
\newcommand{\cW}{\mathcal{W}}
\newcommand{\cG}{\mathcal{G}}
\newcommand{\cF}{\mathcal{F}}
\newcommand{\PP}{\mathcal{P}}
\newcommand{\barPP}{\overline{\mathcal{P}}}


\title{The 3-Qubit Obstruction Ladder: Hyperbolic Embedding,\\
       Mermin Pentagram, and Two Klein Quartic Bridges\\
       {\large (Paper IX)}}


\author{Zhou-Li Chen \\ co-nlang Research}
\date{May 2026}

\begin{document}

\maketitle

\begin{abstract}
  Paper~VIII~\cite{PaperVIII} established the 2-qubit obstruction ladder via
  the explicit transgression functor $\Phi: H^1(K_{3,3}, \Z/2) \to H^2(\CP^3, \Z/2)$.
  The connection to modular geometry was left open: $W(3, \F_2)$ has automorphism
  group $PSp(4, \F_2) \cong S_6$, which is too small to contain $PSL(2, 7)$
  (the Klein quartic's automorphism group).

  Paper~IX steps up to 3 qubits. The hyperbolic embedding
  $GL(3, \F_2) \hookrightarrow PSp(6, \F_2)$ via $g \mapsto \mathrm{diag}(g, g^{-T})$
  (Theorem~\ref{thm:hyperbolic}) realizes the Fano plane $PG(2, \F_2)$ as an
  isotropic subgeometry of $W(5, \F_2)$ (Theorem~\ref{thm:fano-isotropic}).
  Since $GL(3, \F_2) \cong PSL(2, 7)$, this provides a bridge from the 3-qubit
  symplectic geometry to the Klein quartic $X(7)$.

  The Mermin pentagram embeds into $W(5, \F_2)$ with 12{,}096 distinct
  copies, and the $GL(3, \F_2)$-orbit of the GHZ context yields a
  surjective map from pentagram contexts to Fano points
  (Proposition~\ref{prop:fano-pentagram}).

  A second path to $PSL(2, 7)$ arises through the Split Cayley Hexagon
  $H(2)$: $PSL(2, 7) \leq PSU(3, 3) = G_2(2)' \leq G_2(2) \leq PSp(6, \F_2)$.
  We prove (Proposition~\ref{prop:non-conjugacy}) that the two $PSL(2, 7)$
  subgroups are \emph{not} conjugate in $PSp(6, \F_2)$: Path~A stabilizes a
  Lagrangian (order-count argument: $|G_2(2)| = 12{,}096 >
  |P_L| = 10{,}752$ for every Lagrangian stabilizer $P_L$, so Path~B's
  $PSL(2,7) \leq G_2(2)$ cannot stabilize any Lagrangian). Thus
  $PSp(6, \F_2)$ contains \emph{two distinct} Klein quartic bridges.

  The 3-qubit obstruction class $[f_3]$ is formulated in sheaf cohomology
  $H^1(K_5, \cF)$ where $\cF$ is the sheaf of valid eigenvalue assignments
  on the Mermin pentagram nerve.
\end{abstract}

%------------------------------------------------------------------
\section{Introduction}
\label{sec:intro}
%------------------------------------------------------------------

Paper~VIII~\cite{PaperVIII} completed the explicit construction of the
functor $\Phi$ mapping the Peres--Mermin~\cite{Peres1991} contextuality obstruction to
the mod-2 Chern class of twistor space. The 2-qubit MASA poset is
isomorphic to the symplectic polar space $W(3, \F_2)$ with automorphism
group $PSp(4, \F_2) \cong S_6$ of order 720.

The Klein quartic $X(7)$, the modular curve of level~7, has
automorphism group $PSL(2, 7) \cong GL(3, \F_2)$ of order 168. Since
$168 \nmid 720$, no embedding $PSL(2, 7) \hookrightarrow S_6$ exists.
The connection to modular geometry therefore \emph{only opens at 3 qubits},
where $PSp(6, \F_2)$ has order 1{,}451{,}520.

Paper~IX resolves this connection via the \emph{hyperbolic embedding}:
\[
GL(3, \F_2) \hookrightarrow PSp(6, \F_2), \qquad
g \mapsto \begin{pmatrix} g & 0 \\ 0 & g^{-T} \end{pmatrix}.
\]
The main results are:

\begin{enumerate}
  \item \textbf{Theorem~\ref{thm:hyperbolic} (Hyperbolic Embedding).}
    The map $g \mapsto \mathrm{diag}(g, g^{-T})$ is an explicit embedding
    $GL(3, \F_2) \hookrightarrow PSp(6, \F_2)$ preserving the natural
    symplectic form on $\F_2^3 \oplus (\F_2^3)^*$.

  \item \textbf{Theorem~\ref{thm:fano-isotropic} (Fano Isotropic Subgeometry).}
    Under the hyperbolic embedding, the 7 points of the Fano plane
    $PG(2, \F_2)$ embed as an isotropic subgeometry of $W(5, \F_2)$.
    The stabilizer of this subgeometry is $GL(3, \F_2) \cong PSL(2, 7)$.

  \item \textbf{Proposition~\ref{prop:fano-pentagram} (Fano--Pentagram Correspondence).}
    Every Mermin pentagram context in the $GL(3, \F_2)$-orbit of the GHZ
    context contains exactly one point of the Fano isotropic subgeometry.
    The equivariant map from orbit contexts to Fano points is surjective.

  \item \textbf{Proposition~\ref{prop:non-conjugacy} (Non-conjugacy of Two Paths).}
    Path~A's $PSL(2,7)$ (stabilizes the Lagrangian $V$) and Path~B's
    $PSL(2,7)$ (embedded in $G_2(2)$, stabilizes no Lagrangian since
    $|G_2(2)| > |P_L|$) are \emph{not} conjugate in $PSp(6, \F_2)$.
    $PSp(6, \F_2)$ contains \emph{two distinct} Klein quartic bridges.

  \item \textbf{Mock modular forms as obstruction signatures} (\S\ref{sec:mock}):
    The shadow map of mock theta functions is structurally analogous to
    the $d_1$ differential in the Leray spectral sequence.

  \item \textbf{Quantum Darwinism categorification} (\S\ref{sec:darwinism}):
    The adjunction $\mathcal{Q} \dashv \mathcal{B}$ is interpreted as
    the categorical formulation of quantum Darwinism.
\end{enumerate}

The core narrative of this paper:

\begin{quote}
Paper~VIII established the 2-qubit obstruction ladder. Paper~IX steps up
to 3 qubits: the symplectic polar space $W(5, \F_2)$ contains the Fano
plane as an isotropic subgeometry (Theorems~\ref{thm:hyperbolic}--\ref{thm:fano-isotropic}),
whose automorphism group $GL(3, \F_2) \cong PSL(2, 7)$ provides the first
bridge to the Klein quartic. Surprisingly, this ``hyperbolic bridge'' (stabilizes the Lagrangian $V$)
is not conjugate to a second bridge through $G_2(2)$ (which stabilizes
no Lagrangian, since $|G_2(2)| = 12{,}096 > |P_L| = 10{,}752$) ---
$PSp(6, \F_2)$ contains two geometrically inequivalent Klein quartic
structures (Proposition~\ref{prop:non-conjugacy}).
\end{quote}

%------------------------------------------------------------------
\section{The Hyperbolic Embedding Theorem}
\label{sec:hyperbolic}
%------------------------------------------------------------------

\subsection{Symplectic Polar Spaces and Their Automorphisms}
\label{sec:sp-table}

The symplectic polar space $W(2n-1, \F_2)$ is the geometry of totally
isotropic subspaces of $\F_2^{2n}$ equipped with a non-degenerate
symplectic form. Its automorphism group is $PSp(2n, \F_2)$.

\begin{table}[h]
\centering
\begin{tabular}{c l c c l c}
\toprule
$n$ & Space & Points & Lines & Aut Group & Order \\
\midrule
1 & $W(1, \F_2)$ & 3 & --- & $PSp(2, \F_2) \cong S_3$ & 6 \\
2 & $W(3, \F_2) = GQ(2,2)$ & 15 & 15 & $PSp(4, \F_2) \cong S_6$ & 720 \\
3 & $W(5, \F_2)$ & 63 & 315 & $PSp(6, \F_2)$ & 1,451,520 \\
\bottomrule
\end{tabular}
\caption{Symplectic polar spaces $W(2n-1, \F_2)$ for small $n$.}
\label{tab:sp-table}
\end{table}

The 2-qubit case $W(3, \F_2)$ was the setting of Papers~VII and
VIII~\cite{PaperVII, PaperVIII}. The 3-qubit case $W(5, \F_2)$ is the
new territory.

\subsection{The Hyperbolic Embedding}
\label{sec:hyperbolic-embedding}

Let $V = \F_2^3$. The natural symplectic form on $V \oplus V^*$ is:
\[
\langle (v, \phi), (w, \psi) \rangle = \phi(w) + \psi(v) \in \F_2.
\]

The group $GL(3, \F_2)$ acts on $V \oplus V^*$ by:
\[
g \cdot (v, \phi) = (gv, \phi \circ g^{-1}).
\]

This action preserves the symplectic form (direct verification). In
matrix form:
\[
g \mapsto \begin{pmatrix} g & 0 \\ 0 & g^{-T} \end{pmatrix} \in Sp(6, \F_2).
\]

\begin{theorem}[Hyperbolic Embedding]
\label{thm:hyperbolic}
The map $g \mapsto \mathrm{diag}(g, g^{-T})$ is an injective homomorphism
$GL(3, \F_2) \hookrightarrow Sp(6, \F_2) = PSp(6, \F_2)$.
\end{theorem}

\begin{proof}
Injectivity is clear since $g \mapsto g$ is injective on $GL(3, \F_2)$.
The symplectic preservation is verified by:
\[
\langle g(v, \phi), g(w, \psi) \rangle
= (\phi \circ g^{-1})(gw) + (\psi \circ g^{-1})(gv)
= \phi(w) + \psi(v)
= \langle (v, \phi), (w, \psi) \rangle. \qedhere
\]
\end{proof}

\subsection{Fano Plane as Isotropic Subgeometry}
\label{sec:fano-isotropic}

The Fano plane $PG(2, \F_2)$ is the projective plane over $\F_2$:
7 points (non-zero vectors of $\F_2^3$ modulo scalars) and 7 lines
(2-dimensional subspaces).

Under the hyperbolic embedding, the 7 points of $PG(2, \F_2)$ map to
the 7 non-zero vectors of $V \subset V \oplus V^*$. Since $V$ is a
maximal isotropic subspace of $V \oplus V^*$ (the symplectic form
vanishes identically on $V$), the image of $PG(2, \F_2)$ is an
\emph{isotropic subgeometry} of $W(5, \F_2)$.

\begin{theorem}[Fano Isotropic Subgeometry]
\label{thm:fano-isotropic}
Under the hyperbolic embedding $GL(3, \F_2) \hookrightarrow PSp(6, \F_2)$,
the Fano plane $PG(2, \F_2)$ embeds as an isotropic subgeometry of
$W(5, \F_2)$. The stabilizer of this subgeometry in $PSp(6, \F_2)$ is
isomorphic to $GL(3, \F_2) \cong PSL(2, 7)$.
\end{theorem}

\begin{proof}
$V \subset V \oplus V^*$ is maximal isotropic: for all $v, w \in V$,
$\langle (v, 0), (w, 0) \rangle = 0 + 0 = 0$. The projectivization
$\mathbb{P}(V) \cong PG(2, \F_2)$ is therefore an isotropic subgeometry.

The stabilizer in $PSp(6, \F_2)$ is the parabolic subgroup $P$ preserving
the flag $0 \subset V \subset V \oplus V^*$, with Levi decomposition
$P = L \ltimes U$ where $L \cong GL(3, \F_2)$ and $U$ is the unipotent
radical (matrices $\begin{psmallmatrix} I & B \\ 0 & I \end{psmallmatrix}$
with $B = B^T$). Elements of $U$ send $(v, 0)$ to $(v, Bv)$; this vector
remains isotropic (since $\langle (v,Bv), (w,Bw) \rangle = v^T(B+B^T)w = 0$
for symmetric $B$), but it lies outside $V$ unless $B = 0$. Thus $U$
preserves the isotropic property but moves $V$ to a different Lagrangian
subspace. The stabilizer of $\mathbb{P}(V)$ itself (not just any
Lagrangian) is therefore exactly $L \cong GL(3, \F_2)$. \qedhere
\end{proof}

\subsection{The Klein Quartic Connection}
\label{sec:klein-connection}

The Klein quartic $X(7)$ is the modular curve of level 7:
\[
x^3 y + y^3 z + z^3 x = 0 \subset \CP^2.
\]

It is a Riemann surface of genus 3 with automorphism group $PSL(2, 7)$
of order 168. The isomorphism $PSL(2, 7) \cong GL(3, \F_2)$ is classical
(both are simple groups of order 168).

The geometric chain is now complete:
\[
W(5, \F_2) \xrightarrow{\text{Aut}} PSp(6, \F_2)
\supset GL(3, \F_2) \cong PSL(2, 7) \cong \text{Aut}(X(7)).
\]

The 3-qubit MASA poset (63 points, 315 lines in $W(5, \F_2)$) contains
the Fano plane (7 points, 7 lines) as an isotropic subgeometry, and the
automorphism group of this subgeometry is the same as the automorphism
group of the Klein quartic.

\subsection{Relation to 2-qubit $W(3, \F_2)$}
\label{sec:2qubit-relation}

The 2-qubit case $W(3, \F_2)$ has automorphism group
$PSp(4, \F_2) \cong S_6$ of order 720. The relation to $PSL(2, 7)$
(order 168) is indirect:

\begin{itemize}
  \item $720 \div 168 = 4.285\ldots$, not an integer. By Lagrange's
    theorem, $PSL(2, 7) \not\leq S_6$.
  \item $S_6$ contains no subgroup of order 168. The maximal subgroups
    of $S_6$ are $S_5$ (120), $A_6$ (360), and $S_4 \times S_2$ (48) ---
    none divisible by 168.
\end{itemize}

This explains why the Klein quartic connection \emph{only opens at
3-qubit}: the 2-qubit automorphism group $S_6$ is too small to contain
$PSL(2, 7)$.

A computational observation: if one attempts to embed the Fano plane
$PG(2, \F_2)$ into $W(3, \F_2)$ as a point set, only 3 of the 7 Fano
lines are isotropic lines of $W(3, \F_2)$. The remaining 4 Fano lines
fail the isotropic condition. This confirms that the Fano plane's
``natural home'' is $W(5, \F_2)$: the hyperbolic embedding of
Theorem~\ref{thm:hyperbolic} sends all 7 Fano lines to isotropic lines,
which $W(3, \F_2)$ cannot accommodate.

\subsection{Mermin Pentagram and the Fano Subgeometry}
\label{sec:mermin-fano}

The Mermin pentagram (10 operators, 5 contexts of 4) embeds perfectly
into $W(5, \F_2)$. There are exactly \textbf{12{,}096} distinct
pentagram copies in the 3-qubit Pauli group~\cite{LevayPlanatSaniga}.
This number is significant:
\[
12096 = |G_2(2)| = 2^6 \cdot 3^3 \cdot 7.
\]

The Split Cayley Hexagon $H(2)$ is a subgeometry of $W(5, \F_2)$ with
63 points (all of them) and 63 lines (a subset of the 315 lines of
$W(5, \F_2)$). Its automorphism group is $G_2(2)$, and the 12096
pentagrams form a $G_2(2)$-torsor.

The GHZ context $\{XXX, XYY, YXY, YYX\}$ provides a concrete link to
the Fano subgeometry. In the symplectic decomposition
$\F_2^3 \oplus (\F_2^3)^*$:

\[
\begin{array}{c|c|c|c}
\text{Operator} & v \text{ (X part)} & \phi \text{ (Z part)} & \text{In Fano?} \\ \hline
XXX & (1,1,1) & (0,0,0) & \checkmark \\
XYY & (1,1,1) & (0,1,1) & \times \\
YXY & (1,1,1) & (1,0,1) & \times \\
YYX & (1,1,1) & (1,1,0) & \times
\end{array}
\]

All four operators share the same $v$-component $(1,1,1)$. Only $XXX$
has $\phi = 0$, placing it on the Fano plane. Under the $GL(3, \F_2)$
action $g \cdot (v, \phi) = (gv, g^{-T}\phi)$, the orbit of the GHZ
context has the property that each context contains exactly one Fano point.

\begin{proposition}[Fano--Pentagram Correspondence]
\label{prop:fano-pentagram}
Under the hyperbolic embedding, every Mermin pentagram context in the
$GL(3, \F_2)$-orbit of the GHZ context contains exactly one point of
the Fano isotropic subgeometry. The equivariant map
\[
\{\text{orbit contexts}\} \longrightarrow PG(2, \F_2), \quad
C \mapsto C \cap \mathcal{F}
\]
is surjective onto all 7 Fano points.
\end{proposition}

\begin{proof}
The GHZ context $\{XXX, XYY, YXY, YYX\}$ has all four operators sharing
the same $v$-component $(1,1,1)$, with $Z$-components
$(0,0,0), (0,1,1), (1,0,1), (1,1,0)$ respectively. Under the action
$g \cdot (v, \phi) = (gv, g^{-T}\phi)$, the $v$-components of all four
operators in $g \cdot C_{\text{GHZ}}$ are identically $g(1,1,1)$, while
the $Z$-components are $g^{-T}\phi_i$ for the four distinct $\phi_i$.
Since $g^{-T}$ is a bijection, the four $Z$-components remain distinct;
exactly one is zero (the image of $\phi = 0$). The Fano plane condition
is $\phi = 0$, so each context contains exactly one Fano point.

Surjectivity follows because $GL(3, \F_2)$ acts transitively on the
7 non-zero vectors of $\F_2^3$, hence transitively on the 7 Fano points.
\qedhere
\end{proof}

\begin{remark}[Pentagram--Fano Nerve: Toward a Finite Geometry Dictionary]
\label{rem:nerve-fano-open}
Proposition~\ref{prop:fano-pentagram} maps \emph{GHZ-orbit} contexts
to all 7 Fano points surjectively. A finer question arises for a
specific Mermin pentagram (only 5 contexts): which 5 of the 7 Fano
points appear, and what geometric configuration do they form in
$PG(2, \F_2)$?

The $K_5$ nerve has exactly $\binom{5}{2} = 10$ edges (shared operators),
while 5 Fano points determine 10 pairs, each spanning a projective line.
This numerical coincidence suggests a \emph{finite geometry dictionary}:
\[
\begin{aligned}
\{K_5 \text{ edges (shared operators)}\}
&\;\longleftrightarrow\;
\{\text{pairs of Fano points}\} \\
&\;\longleftrightarrow\;
\{\text{incidence lines in } PG(2, \F_2)\}.
\end{aligned}
\]
In $PG(2, \F_2)$, cross-ratios degenerate (characteristic 2), so the
relevant invariant is \emph{collinearity}. Any 5 Fano points contain at
least one collinear triple; these triples partition the 10 pairs into
``collinear'' and ``non-collinear'' classes.

\begin{conjecture}[Collinearity--Obstruction Dictionary]
\label{conj:collinearity-obstruction}
Under the hyperbolic embedding, the non-triviality of
$[f_3] \in H^1(K_5, \cF)$ is equivalent to the collinearity
pattern of the 5 corresponding Fano points admitting no consistent
$\Z/2$-labeling compatible with all incidence relations. The
Kochen--Specker contradiction ($0 = 1 \pmod 2$ from product parity)
translates to a collinearity obstruction in $PG(2, \F_2)$.
\end{conjecture}

If confirmed, this conjecture would yield a purely geometric proof that
$[f_3] \neq 0$, bypassing the algebraic product-parity argument, and
connect the pentagram obstruction directly to the projective geometry of
the Klein quartic's automorphism group $PSL(2, 7) = \mathrm{Aut}(PG(2,\F_2))$.
\end{remark}

\subsection{Split Cayley Hexagon and $G_2(2)$}
\label{sec:cayley-hexagon}

The number 12096 has a deeper geometric meaning. The Split Cayley
Hexagon $H(2)$ is a generalized hexagon embedded in $W(5, \F_2)$:
\begin{itemize}
  \item \textbf{Points:} 63 (all points of $W(5, \F_2)$)
  \item \textbf{Lines:} 63 (a distinguished subset of the 315 lines)
  \item \textbf{Automorphism group:} $G_2(2)$ of order 12096
\end{itemize}

The derived subgroup $G_2(2)' = PSU(3, 3)$ has order 6048. From the
ATLAS of Finite Groups~\cite{ATLAS}, $PSU(3, 3)$ contains $PSL(2, 7)$ as a maximal
subgroup (index 36). This provides a second path from $PSL(2, 7)$ to
$W(5, \F_2)$:
\[
PSL(2, 7) \leq PSU(3, 3) = G_2(2)' \leq G_2(2) \leq PSp(6, \F_2).
\]

\subsection{Two Paths to the Klein Quartic}
\label{sec:two-paths}

Two independent embeddings of $PSL(2, 7) \cong GL(3, \F_2)$ into
$PSp(6, \F_2)$:

\begin{itemize}
  \item \textbf{Path A (hyperbolic):}
    $GL(3, \F_2) \hookrightarrow PSp(6, \F_2)$ via
    $g \mapsto \mathrm{diag}(g, g^{-T})$ (Theorem~\ref{thm:hyperbolic})
  \item \textbf{Path B (Cayley Hexagon):}
    $PSL(2, 7) \leq PSU(3, 3) \leq G_2(2) \leq PSp(6, \F_2)$
\end{itemize}

\begin{proposition}[Non-conjugacy of Paths]
\label{prop:non-conjugacy}
The two $PSL(2, 7)$ subgroups from Path A (hyperbolic) and Path B
(Cayley Hexagon) are \emph{not} conjugate in $PSp(6, \F_2)$. Path~A's
$PSL(2,7)$ stabilizes a Lagrangian (namely $V$); Path~B's $PSL(2,7)$,
embedded via $G_2(2)$, does not stabilize any Lagrangian. Since
$|G_2(2)| = 12{,}096 > |P_L| = 10{,}752$ for every Lagrangian
stabilizer $P_L$, conjugation in $PSp(6, \F_2)$ cannot map one
subgroup to the other.
\end{proposition}

\begin{proof}
\textbf{Path A fixes a Lagrangian; Path B cannot.}

The hyperbolic embedding $g \mapsto \mathrm{diag}(g, g^{-T})$ preserves
the Lagrangian subspace $V \subset V \oplus V^*$ by construction, so
Path~A's $PSL(2,7)$ is contained in the parabolic subgroup
$P_V \leq Sp(6, \F_2)$ stabilizing $V$.

The group $Sp(6, \F_2)$ acts transitively on the 135 Lagrangians of
$W(5, \F_2)$, giving $|P_V| = |Sp(6, \F_2)|/135 = 1{,}451{,}520/135 = 10{,}752$.
Since $|G_2(2)| = 12{,}096 > 10{,}752$, the subgroup $G_2(2)$ is strictly
larger than any Lagrangian stabilizer: $G_2(2) \not\leq P_L$ for
\emph{every} Lagrangian $L$. In particular, $G_2(2)$ fixes no Lagrangian.

Path~B's $PSL(2,7)$ is embedded in $G_2(2)$, whose geometric structure
(automorphisms of the Split Cayley Hexagon) is incompatible with fixing
any Lagrangian. Path~B therefore does not stabilize $V$ (nor any other
Lagrangian in the same $PSp(6,\F_2)$-orbit as $V$).

Since Path~A's $PSL(2,7)$ stabilizes a Lagrangian and Path~B's does
not, the two subgroups cannot be conjugate in $PSp(6, \F_2)$: conjugation
by any element of $PSp(6, \F_2)$ maps stabilizers of Lagrangians to
stabilizers of Lagrangians. \qedhere
\end{proof}

This means $PSp(6, \F_2)$ contains \emph{two distinct} Klein quartic
bridges, each with its own geometric significance:

\begin{table}[h]
\centering
\begin{tabular}{l l l}
\toprule
& Path A (hyperbolic) & Path B ($G_2(2)$) \\
\midrule
Representation & Reducible $3 \oplus 3^*$ & 6-dim, no invariant Lagrangian \\
Fixed geometry & Fano plane $\mathbb{P}(V)$ & No invariant Lagrangian \\
Semantics & Classical $PSL(2,7)$ & Quantum $PSL(2,7)$ \\
\bottomrule
\end{tabular}
\caption{Two non-conjugate Klein quartic bridges in $PSp(6, \F_2)$.}
\label{tab:two-paths}
\end{table}

\begin{itemize}
  \item \textbf{Path A} (hyperbolic): connects to the Fano isotropic
    subgeometry via the Lagrangian decomposition
  \item \textbf{Path B} (Cayley Hexagon): connects to the Split Cayley
    Hexagon $H(2)$ and its 12096 pentagrams
\end{itemize}

\begin{remark}[Path B and Non-Geometric Backgrounds ($\circ$)]
\label{rem:pathb-physics}
Path~A's Lagrangian-stabilizing $PSL(2,7)$ has a natural physical
interpretation: the decomposition $\F_2^6 = V \oplus V^*$ mirrors the
position/momentum split in T-duality, and stabilizing $V$ corresponds
to fixing a ``geometric'' D-brane boundary condition.

Path~B, embedded via $G_2(2)$, stabilizes no Lagrangian. In string
theory, the failure to preserve any Lagrangian under T-duality is
precisely the signature of \emph{non-geometric fluxes}. In the chain
$H \to f \to Q \to R$, the $Q$- and $R$-fluxes are those that mix
position and momentum coordinates in a way incompatible with any
geometric T-duality frame. Path~B's $PSL(2,7) \leq G_2(2)$ may
therefore correspond to the discrete symmetry of a non-geometric flux
background --- the two Klein quartic bridges would represent:
Path~A $=$ geometric modular geometry, Path~B $=$ its non-geometric
flux dual. Whether $G_2(2)$ (automorphisms of the Split Cayley Hexagon)
appears as a discrete symmetry in M-theory compactification on a
$G_2$-holonomy manifold is an open problem (\S\ref{sec:pathb-physics}).
\end{remark}

%------------------------------------------------------------------
\section{Mock Modular Forms as $H^1$ Obstruction Signatures}
\label{sec:mock}
%------------------------------------------------------------------

\subsection{Mock Theta Functions and the Shadow Map}

A mock theta function $f(\tau)$ is a holomorphic function on the upper
half-plane that ``almost'' transforms like a modular form but requires
a non-holomorphic correction term~\cite{Zwegers}. The completion
$\widehat{f} = f + g^*$ transforms as a modular form, where $g$ is the
\emph{shadow} --- a genuine modular form of weight $k + 1/2$.

The shadow map is a linear operator:
\[
\xi: \widehat{M}_k \longrightarrow S_{k+1/2},
\]
where $\widehat{M}_k$ is the space of completed mock modular forms and
$S_{k+1/2}$ is the space of cusp forms.

The exact sequence:
\[
0 \longrightarrow M^!_k \longrightarrow \widehat{M}_k
\xrightarrow{\;\xi\;} S_{k+1/2} \longrightarrow 0
\]
has the structure of an $H^1$ exact sequence: mock forms are 1-cocycles
that are not coboundaries, and the shadow measures ``how far'' they are
from being genuine modular forms.

\subsection{Shadow Map as $d_1$ Differential}

The structural parallel with the obstruction ladder:

\begin{table}[h]
\centering
\begin{tabular}{l l}
\toprule
Mock Modular Forms & Obstruction Ladder \\
\midrule
Shadow map $\xi$ & $d_1$ differential (Leray SS) \\
$\widehat{f}$ ``non-holomorphic'' & Gerbe-twisted $H^1$ non-vanishing \\
Level 7 modular equations & Klein quartic $X(7)$ with $PSL(2, 7)$ action \\
$p(n) \bmod 2$ distribution & $H^2(\CP^3, \Z/2)$ class \\
Rogers--Ramanujan identities & Self-dual / anti-self-dual splitting \\
\bottomrule
\end{tabular}
\caption{Structural parallel between mock modular forms and the obstruction ladder.}
\label{tab:mock-parallel}
\end{table}

The shadow map $\xi$ is structurally identical to the $d_1$ differential
in the Leray spectral sequence: both measure the failure of local data
to glue into global sections.

\subsection{Level 7 and the Klein Quartic}

The modular curve $X(7)$ is the Klein quartic. The space of cusp forms
$S_{3/2}(\Gamma(7))$ carries an action of $PSL(2, 7)$. The shadow map
for level-7 mock forms:
\[
\xi_7: \widehat{M}_k(\Gamma(7)) \longrightarrow S_{k+1/2}(\Gamma(7))
\]
is equivariant under the $PSL(2, 7)$ action.

\begin{remark}[Shadow--Obstruction Correspondence]
\label{rem:shadow-obstruction}
The shadow map $\xi_7$ for level-7 mock modular forms is structurally
analogous to the $d_1$ differential in the Leray spectral sequence for
the Fano isotropic subgeometry of $W(5, \F_2)$:
\[
\xi_7 \;\text{``}\cong\text{''}\; d_1: E_1^{0,1} \longrightarrow E_1^{1,1}.
\]
This is a heuristic correspondence: $\xi_7$ is a linear map between
infinite-dimensional complex function spaces (half-integral weight
modular forms), while $d_1$ is a differential between finite
$\Z/2$-vector spaces. An equivariant isomorphism would require a
precisely defined functor from modular forms to $\F_2$-cohomology
(e.g., a mod-$p$ reduction functor), which is not yet constructed.
\end{remark}

\subsection{Partition Function: Open Question}

Bringmann and Ono~\cite{BringmannOno2006} proved the Andrews--Dragonette conjecture using
traces of singular moduli. Singular moduli are elements of the ideal
class group --- which is an $H^1$ obstruction group.

Ramanujan's celebrated congruences establish deep modular structure for
$p(n)$: $p(5n+4) \equiv 0 \pmod 5$, $p(7n+5) \equiv 0 \pmod 7$,
$p(11n+6) \equiv 0 \pmod{11}$. The behavior modulo 2 is more subtle:
$p(n)$ is odd for a density-zero set of integers (a theorem of Serre
via the theory of lacunary modular forms). The critical alignment for
this paper is that $H^2(\CP^3, \Z/2)$ is also $\Z/2$-valued: any
pairing $H_1(K_{3,3}, \Z/2) \times H^1(K_{3,3}, \Z/2) \to \Z/2$
produces \emph{parity} outputs, naturally matching the $p(n) \bmod 2$
domain.

A speculative connection: the partition function $p(n)$ modulo 2 might
be related to the obstruction class of the 2-qubit MASA poset via a
pairing:
\[
p(n) \bmod 2 \stackrel{?}{=} \langle [f], \text{cycle}_n \rangle_{\Z/2}.
\]
However, this requires an explicit construction of
$\text{cycle}_n \in H_1(K_{3,3}, \Z/2)$ from the binary expansion of
$n$. This is moved to \S\ref{sec:open-problems} as an open problem.

%------------------------------------------------------------------
\section{$\mathcal{Q} \dashv \mathcal{B}$ as Quantum Darwinism Categorification}
\label{sec:darwinism}
%------------------------------------------------------------------

\subsection{The Adjunction}

Paper~V~\cite{PaperV} conjectured an adjunction $\mathcal{Q} \dashv \mathcal{B}$
between ``quantum'' and ``classical'' categories. Paper~VIII realized
this for the $SU(2)$ case:
\begin{itemize}
  \item $\mathcal{Q} = \Phi_*$: quantization (pushforward from MASA nerve
    to twistor space)
  \item $\mathcal{B} = \Phi^*$: Bohrification (pullback from quantum to
    classical)
\end{itemize}

\subsection{Quantum Darwinism}

Quantum Darwinism (Zurek~\cite{Zurek2009}) asks: what information survives decoherence
and becomes accessible to multiple observers? The core mechanism is
\emph{redundant encoding} --- classical information is the part of the
quantum state that can be reconstructed from many disjoint environment
fragments.

The adjunction $\mathcal{Q} \dashv \mathcal{B}$ categorifies this:
\begin{itemize}
  \item $\mathcal{B} = \Phi^*$ (Bohrification) = \emph{coarse-graining}:
    projects the full quantum state onto classical measurement contexts
    (MASAs)
  \item $\mathcal{Q} = \Phi_*$ (quantization) = \emph{record lifting}:
    attempts to reconstruct the quantum state from classical records
\end{itemize}

\subsection{PM Obstruction as Non-extendable Section}

The PM obstruction class $[f] \in H^1(K_{3,3}, \Z/2)$ is precisely the
information that \emph{cannot} be reconstructed from local MASA records.
In the language of Quantum Darwinism:

\begin{itemize}
  \item Each MASA context is an ``environment fragment''
  \item The sheaf of local sections over the MASA nerve represents the
    classical records
  \item The obstruction $[f]$ is the \emph{non-extendable section} ---
    the part of the quantum state that cannot be redundantly encoded
\end{itemize}

\begin{proposition}[Darwinistic Obstruction]
\label{prop:darwinism}
The PM contextuality class $[f] \in H^1(K_{3,3}, \Z/2)$ measures the
information that is lost under Bohrification and cannot be recovered by
any classical record-lifting procedure.
\end{proposition}

\begin{proof}
The sheaf of local sections over $K_{3,3}$ has no global section
extending the local data on the 6-cycle. The obstruction class $[f]$
is a specific non-zero element of
$H^1(K_{3,3}, \Z/2) \cong (\Z/2)^4$ --- the PM obstruction class ---
that prevents this extension (Lemma~2.3 of Paper~VIII~\cite{PaperVIII}).
\qedhere
\end{proof}

\subsection{n-Qubit Generalization}

For $n$ qubits, the MASA nerve is the \v{C}ech nerve of
$W(2n-1, \F_2)$. The obstruction class
$[f_n] \in H^1(\mathcal{N}(\mathcal{C}(\mathcal{P}_n)), \Z/2)$ measures
the information lost under $n$-qubit Bohrification.

For $n=3$: $W(5, \F_2)$ has 63 points and 315 lines. The Fano isotropic
subgeometry (7 points in $V$) is a ``classical core'' that can be
redundantly encoded. The choice of $V$ over $V^*$ is not arbitrary: $V$
corresponds to the position basis (the MASA of diagonal operators in the
computational basis), while $V^*$ corresponds to the momentum basis. In
the Bohrification picture, $V$ is the privileged classical context
selected by the measurement apparatus --- the ``pointer basis'' in
Zurek's terminology. The remaining 56 points of
$W(5, \F_2) \setminus (V \cup V^*)$ carry the quantum obstruction.

%------------------------------------------------------------------
\section{Open Problems}
\label{sec:open-problems}
%------------------------------------------------------------------

\subsection{Gerbe $H^1$ Non-vanishing (from Paper~VIII)}

The twisted Penrose transform
$H^1_{\mathcal{G}}(\CP^3, \cO_{\mathcal{G}}(-n-2))$ remains an open
problem. The standard Leray analysis gives $R^1\pi_*\cO_{\mathcal{G}} = 0$,
so the $\Z/2$-gerbe obstruction cannot be detected by classical sheaf
cohomology. Three concrete activation mechanisms:

\begin{enumerate}
  \item \textbf{Kapustin--Witten $B$-field twist.}
    {\sloppy In the geometric Langlands program~\cite{KapustinWitten2007},
    a $B$-field background (a $U(1)$-gerbe) twists the category
    of $\mathcal{D}$-modules; a $\Z/2$-gerbe corresponds to an
    orientifold $B$-field. In the KW topological twist of
    $\mathcal{N}=4$ SYM, the gerbe class $[e] \in H^2(\CP^3, \Z/2)$
    deforms the de~Rham complex and can activate cohomology that
    vanishes without the twist. Whether $[e]$ gives non-zero
    $H^1_{\mathcal{G}}$ after the KW twist is the key computation.\par}

  \item \textbf{Derived category with singular support.}
    The classical Leray failure ($R^1\pi_*\cO_{\mathcal{G}} = 0$) is an
    ordinary derived category statement. The theory of singular support
    in $D^b_{\mathcal{G}}(\CP^3)$ --- in the spirit of the Arinkin--Gaitsgory
    spectral decomposition in geometric Langlands --- provides a mechanism:
    $R^1\pi_*\cO_{\mathcal{G}}$ may be non-trivial as an object with
    specified singular support, even though it vanishes in the full
    derived category. The $\Z/2$-Clifford algebra twist is the natural
    coefficient system.

  \item \textbf{Ambitwistor space and orientifold flux.}
    In ambitwistor space $\mathbb{PA} \subset \mathbb{PT} \times \mathbb{PT}^*$,
    the NS-NS $B$-field is a $U(1)$-gerbe; a $\Z/2$-gerbe corresponds to
    orientifold flux. The twisted Penrose transform in this setting may
    yield non-vanishing $H^1_{\mathcal{G}}$ via discrete instanton
    corrections to the flat-space result.
\end{enumerate}

\subsection{3-qubit $W(5, \F_2)$ Obstruction Class}
\label{sec:3qubit-obstruction}

The 3-qubit analog of the PM square is the \emph{Mermin pentagram}
(Mermin~\cite{Mermin1990}): 10 three-qubit Pauli operators arranged on a pentagram
with 5 contexts (lines of 4 mutually commuting operators). Each operator
appears in exactly 2 contexts.

One context is
$\{X \otimes X \otimes X,\; X \otimes Y \otimes Y,\; Y \otimes X \otimes Y,\;
Y \otimes Y \otimes X\}$, whose product is $-\mathbf{I}_8$. In the
symplectic vector space $\F_2^6$, these four operators span a
3-dimensional isotropic subspace (a Lagrangian):
\[
\begin{array}{c|c}
\text{Operator} & (x_1,z_1,x_2,z_2,x_3,z_3) \\ \hline
XXX & (1,0,1,0,1,0) \\
XYY & (1,0,1,1,1,1) \\
YXY & (1,1,1,0,1,1) \\
YYX & (1,1,1,1,1,0)
\end{array}
\]
The sum $v_1 + v_2 + v_3 = v_4 \pmod 2$, confirming linear dependence.
The full MASA is
\[
\{XXX, XYY, YXY, YYX, IZZ, ZIZ, ZZI\}
\]
(7 non-identity elements).

\begin{table}[h]
\centering
\begin{tabular}{c l}
\toprule
Context & 4 Operators \\
\midrule
$C_1$ & $XXX,\; XYY,\; YXY,\; YYX$ \\
$C_2$ & $XXX,\; ZZX,\; ZXZ,\; XZZ$ \\
$C_3$ & $YXY,\; ZXZ,\; ZYY,\; YYZ$ \\
$C_4$ & $YYX,\; ZZX,\; ZYY,\; YZY$ \\
$C_5$ & $XYY,\; XZZ,\; YYZ,\; YZY$ \\
\bottomrule
\end{tabular}
\caption{Explicit Mermin pentagram: 10 operators across 5 contexts.}
\label{tab:pentagram}
\end{table}

Shared operators (= $K_5$ edges):
\begin{align*}
C_1 \cap C_2 &= XXX, & C_1 \cap C_3 &= YXY, &
C_1 \cap C_4 &= YYX, & C_1 \cap C_5 &= XYY, \\
C_2 \cap C_3 &= ZXZ, & C_2 \cap C_4 &= ZZX, &
C_2 \cap C_5 &= XZZ, \\
C_3 \cap C_4 &= ZYY, & C_3 \cap C_5 &= YYZ, &
C_4 \cap C_5 &= YZY.
\end{align*}

All 10 pairs within each context commute (verified: each context's 4
operators share a common eigenbasis). The product of operators in each
context yields $\pm \mathbf{I}_8$; the parity of minus signs encodes the
obstruction class $[f_3]$.

\begin{remark}[KS Contradiction Prerequisite]
{\sloppy The Kochen--Specker contradiction~\cite{KochenSpecker1967} requires
an \emph{odd} number of contexts with product $-\mathbf{I}$. For the
pentagram in Table~\ref{tab:pentagram}, the context products must be
verified: the system of 5 context equations is inconsistent because
summing all 5 equations gives $0 = 1 \pmod 2$ (each of the 10 operators
appears in exactly 2 contexts, so the LHS sums to 0, but the RHS sums
to 1 from the anomalous context(s)). The explicit count of $-\mathbf{I}$
contexts should be odd for the KS contradiction to hold.\par}
\end{remark}

The \v{C}ech nerve of the Mermin pentagram is the complete graph $K_5$
(5 vertices = contexts, 10 edges = shared operators):
\[
H^1(K_5, \Z/2) \cong (\Z/2)^6, \qquad \beta_1 = 6.
\]
The geometric transition from 2-qubit to 3-qubit is $K_{3,3}$
(bipartite, row/column symmetry) $\to K_5$ (complete, five-fold
symmetry).

The obstruction class $[f_3]$ is the cohomology class corresponding to
the $-\mathbf{I}$ product anomaly. Unlike the 2-qubit case where
$[f] \in H^1(K_{3,3}, \Z/2)$ with constant coefficients, the 3-qubit
obstruction lives in the \emph{sheaf cohomology} $H^1(K_5, \cF)$ where
$\cF$ is the sheaf of valid eigenvalue assignments:

\begin{itemize}
  \item For each context $C_i$,
    $\cF(C_i) = \{ \text{valid } \pm 1 \text{ assignments to 4 operators} \}
    \cong (\Z/2)^3$ (4 variables with 1 product constraint).
  \item For each edge $C_i \cap C_j$, $\cF(C_i \cap C_j) = \{ \pm 1 \}$
    (the shared operator's eigenvalue).
  \item Restriction maps: project a context's assignment to the shared
    operator's value.
\end{itemize}

The obstruction class $[f_3] \in H^1(K_5, \cF)$ measures the failure of
local sections to glue into a global section. For the pentagram,
$[f_3] \neq 0$ because the system of 5 context equations is inconsistent:
summing all 5 equations gives $0 = 1 \pmod 2$ (each of the 10 operators
appears in exactly 2 contexts, so the LHS sums to 0, but the RHS sums
to 1 from the single anomalous context).

Under the hyperbolic embedding (Theorem~\ref{thm:hyperbolic}), $[f_3]$
should map to a class in the equivariant cohomology of the Fano isotropic
subgeometry:
\[
[f_3] \in H^1_{GL(3, \F_2)}(W(5, \F_2), \Z/2).
\]
The relation to the Klein quartic
(Remark~\ref{rem:shadow-obstruction}) is the key
open computation.

\subsection{Partition Function Modulo 2}

The speculative connection
$p(n) \bmod 2 = \langle [f], \text{cycle}_n \rangle_{\Z/2}$ requires:
\begin{enumerate}
  \item An explicit construction of $\text{cycle}_n \in H_1(K_{3,3}, \Z/2)$
    from the binary expansion of $n$
  \item Verification that the pairing output matches the distribution of
    $p(n) \bmod 2$ (an active research problem in number theory; the
    known result is that $p(n)$ is odd for a density-zero set, via
    Serre's theory of lacunary modular forms)
\end{enumerate}

The motivation for $\bmod 2$ specifically is the parity alignment:
$H^2(\CP^3, \Z/2)$ is $\Z/2$-valued, and the obstruction pairing
lands in $\Z/2$. Ramanujan's congruences establish that $p(n)$ has
non-trivial modular structure at small primes; whether the $\bmod 2$
case connects to the obstruction class via the $K_{3,3}$ nerve is a
wishlist item, not a conjecture --- the mapping from integers to
cycles is not yet defined.

\subsection{n-Qubit Ladder}

The general pattern for the obstruction ladder:
\begin{itemize}
  \item $n=1$: $W(1, \F_2)$, $H^1$, geometric phase (Paper~I)
  \item $n=2$: $W(3, \F_2)$, $H^2$, KS contextuality
    (Papers~III, VII, VIII~\cite{PaperIII, PaperVII, PaperVIII})
  \item $n=3$: $W(5, \F_2)$, $H^3$, Borromean non-associativity
    (Paper~IV~\cite{PaperIV}, open)
  \item $n \geq 4$: $W(2n-1, \F_2)$, higher obstructions (open)
\end{itemize}

The modular geometry connection (via hyperbolic embedding) only opens
at $n=3$, where $W(5, \F_2)$ contains the Fano plane. This suggests
that the Ramanujan/modular connection is intrinsically a \emph{3-qubit
and above} phenomenon.

\subsection{Path B and Non-Geometric Fluxes ($\circ$)}
\label{sec:pathb-physics}

Proposition~\ref{prop:non-conjugacy} establishes two non-conjugate
$PSL(2,7)$ subgroups in $PSp(6, \F_2)$: Path~A stabilizes a Lagrangian,
Path~B (via $G_2(2)$) does not. The open problem is to give Path~B
a concrete physical interpretation.

The T-duality classification of string flux backgrounds suggests a
natural dictionary. A Lagrangian subspace $L \subset \F_2^{2n}$ defines
a ``geometric'' T-duality frame: T-dualizing along directions within $L$
maps geometric compactifications to geometric ones ($H$- and $f$-flux).
When no Lagrangian is preserved under the symmetry group, one enters
the domain of \emph{non-geometric fluxes} ($Q$-flux and $R$-flux),
where the background cannot be described in any single geometric duality
frame.

\begin{itemize}
  \item If Path~B's $PSL(2,7) \leq G_2(2) \leq PSp(6, \F_2)$ is the
    discrete symmetry group of a non-geometric flux background, the two
    Klein quartic bridges would represent \emph{geometric} vs.\
    \emph{non-geometric} modular structures on the same underlying algebra.
  \item The group $G_2(2)$ (automorphisms of the Split Cayley Hexagon)
    might appear as a discrete symmetry in M-theory compactification on
    a $G_2$-holonomy manifold, where the $G_2$ holonomy group and
    $G_2(2)$ (the split real form of $G_2$) are related by real-form
    variation.
  \item The obstruction class $[f_3] \in H^1(K_5, \cF)$ should map to
    the discrete torsion of the non-geometric background if the
    dictionary is valid.
\end{itemize}

Required: a concrete $G_2$-holonomy compactification whose discrete
symmetry is $G_2(2)$ of order 12{,}096, and a proof that the
obstruction class maps to the corresponding discrete torsion group.

%------------------------------------------------------------------
\section*{Appendix: Rigor Self-Assessment}
%------------------------------------------------------------------

\begin{table}[H]
\centering\small
\setlength{\tabcolsep}{4pt}
\begin{tabular}{l c p{6.5cm}}
\toprule
Claim & Rigor & Status \\
\midrule
Hyperbolic embedding $GL(3,\F_2) \hookrightarrow PSp(6,\F_2)$ & $\bullet$ &
  Explicit construction $g \mapsto \mathrm{diag}(g, g^{-T})$; symplectic
  preservation verified \\
Fano plane as isotropic subgeometry of $W(5,\F_2)$ & $\bullet$ &
  $V \subset V \oplus V^*$ is maximal isotropic; projectivization gives
  $PG(2,\F_2)$ \\
$PSL(2,7) \cong GL(3,\F_2)$ & $\bullet$ &
  Classical isomorphism; both simple groups of order 168 \\
Split Cayley Hexagon $H(2) \subset W(5,\F_2)$ & $\bullet$ &
  63 points, 63 lines, $\mathrm{Aut} = G_2(2)$; standard finite geometry \\
Mermin pentagram in $W(5,\F_2)$ (12096 copies) & $\bullet$ &
  $12096 = |G_2(2)|$; each pentagram = 5 contexts of 4 commuting operators \\
Shadow map $\xi$ as $d_1$ differential & $\circ$ &
  Structural parallel established; equivariant isomorphism is a motivating
  analogy, not yet a conjecture \\
$p(n) \bmod 2$ = obstruction pairing & $\circ$ &
  Moved to \S\ref{sec:open-problems} Open Problem; requires explicit cycle
  construction \\
$\mathcal{Q} \dashv \mathcal{B}$ = Quantum Darwinism & $\bullet/\circ$ &
  Adjunction realized for $SU(2)$ ($\bullet$); Darwinism interpretation
  conceptual ($\circ$) \\
PM obstruction = non-extendable section & $\bullet/\circ$ &
  Proposition~\ref{prop:darwinism}: $[f]$ is a specific non-zero class in
  $(\Z/2)^4$ ($\bullet$); Darwinism link conceptual ($\circ$) \\
3-qubit $[f_3]$ construction & $\bullet/\circ$ &
  Mermin pentagram identified (12096 copies); nerve $K_5$ and $H^1$
  structure ($\bullet$); explicit sheaf cohomology description ($\bullet$);
  equivariant class under hyperbolic embedding pending ($\circ$) \\
Proposition~\ref{prop:fano-pentagram} (Fano--Pentagram) & $\bullet$ &
  Proven: each context has exactly one Fano point; surjectivity by
  transitivity \\
Proposition~\ref{prop:non-conjugacy} (Non-conjugacy) & $\bullet$ &
  Proven: Path~A stabilizes Lagrangian $V$; Path~B cannot since
  $|G_2(2)| = 12{,}096 > |P_L| = 10{,}752$; not conjugate \\
Gerbe $H^1$ non-vanishing & $\circ$ &
  Open problem from Paper~VIII \\
$n$-qubit ladder pattern & $\bullet/\circ$ &
  $n=1,2$ verified ($\bullet$); $n \geq 3$ conjectural ($\circ$) \\
Modular connection at $n \geq 3$ & $\circ$ &
  Hyperbolic embedding exists; shadow--obstruction correspondence
  conjectural \\
\bottomrule
\end{tabular}
\caption{Rigor self-assessment for Paper~IX.}
\label{tab:rigor}
\end{table}

%------------------------------------------------------------------
\section*{Acknowledgments}
%------------------------------------------------------------------

The author thanks the mathematical physics reading group for feedback
on the hyperbolic embedding framework, and a research collaborator for
identifying the non-conjugacy of the two Klein quartic bridges and the
sheaf cohomology formulation of the 3-qubit obstruction.

%------------------------------------------------------------------
\begin{thebibliography}{10}
%------------------------------------------------------------------

\bibitem[ATLAS]{ATLAS}
J.~H.~Conway, R.~T.~Curtis, S.~P.~Norton, R.~A.~Parker, and R.~A.~Wilson,
\textit{Atlas of Finite Groups: Maximal Subgroups and Ordinary Characters
for Simple Groups},
Oxford University Press (1985).

\bibitem[BringmannOno2006]{BringmannOno2006}
K.~Bringmann and K.~Ono,
``The $f(q)$ mock theta function conjecture and partition ranks,''
\textit{Invent. Math.} \textbf{165}, 243--266 (2006).

\bibitem[KapustinWitten2007]{KapustinWitten2007}
A.~Kapustin and E.~Witten,
``Electric-Magnetic Duality and the Geometric Langlands Program,''
\textit{Commun.\ Number Theory Phys.} \textbf{1}, 1--236 (2007).

\bibitem[KochenSpecker1967]{KochenSpecker1967}
S.~Kochen and E.~P.~Specker,
``The problem of hidden variables in quantum mechanics,''
\textit{J. Math. Mech.} \textbf{17}, 59--87 (1967).

\bibitem[LevayPlanatSaniga]{LevayPlanatSaniga}
P.~L\'evay, M.~Planat, and M.~Saniga,
``Grassmannian Connection Between Three- and Four-Qubit Observables,
Mermin's Contextuality and Black Holes,''
\textit{JHEP} \textbf{09} (2013) 037.

\bibitem[Mermin1990]{Mermin1990}
N.~D.~Mermin,
``Simple unified form for the major no-hidden-variables theorems,''
\textit{Am. J. Phys.} \textbf{58}, 731--734 (1990).

\bibitem[PaperIII]{PaperIII}
Z.-L.~Chen,
``Kochen--Specker Contextuality as Central Extension:
The Peres--Mermin Square and the Pauli Group,''
\textit{Zenodo, DOI: 10.5281/zenodo.20073127} (2026).

\bibitem[PaperIV]{PaperIV}
Z.-L.~Chen,
``The Cohomological Obstruction Ladder:
Lyndon--Hochschild--Serre Transgression and the $H^3$ Frontier,''
\textit{Zenodo, DOI: 10.5281/zenodo.20073184} (2026).

\bibitem[PaperV]{PaperV}
Z.-L.~Chen,
``Observation as Functor: 
The Adjunction of Quantum and Classical,''
\textit{Zenodo, DOI: 10.5281/zenodo.20073318}, 2026.

\bibitem[PaperVII]{PaperVII}
Z.-L.~Chen,
``Twistor Theory from the Obstruction Ladder:
The Googly Problem as an $H^2$ Obstruction,''
\textit{Zenodo, DOI: 10.5281/zenodo.20438042} (2026).

\bibitem[PaperVIII]{PaperVIII}
Z.-L.~Chen,
``The $\Phi$ Functor and the $n$-Qubit Obstruction Ladder:
Explicit Transgression and Twisted Penrose Transform,''
\textit{Zenodo, DOI: 10.5281/zenodo.20454120} (2026).

\bibitem[Peres1991]{Peres1991}
A.~Peres,
``Incompatible results of quantum measurements,''
\textit{J. Phys. A} \textbf{24}, L175--L178 (1991).

\bibitem[Zurek2009]{Zurek2009}
W.~H.~Zurek,
``Quantum Darwinism,''
\textit{Nat. Phys.} \textbf{5}, 181--188 (2009).

\bibitem[Zwegers]{Zwegers}
S.~Zwegers,
``Mock Theta Functions,''
\textit{Ph.D.\ Thesis, Universiteit Utrecht} (2002).

\end{thebibliography}

\end{document}
